% GENERATED FILE. Edit canonical lessons, then run npm run generate:books.
% canonical-source-hash: fnv1a64:564f5ea55f6630fc
% canonical-lessons: LA-C447-cimex, LA-C447-mordeo, LA-C447-scalpo, LA-C447-insomnis, LA-C447-cubo

\chapter{Bedbug, To bite, To scratch, Sleepless, To lie down}
\label{ch:la-bedbug-to-bite-to-scratch-sleepless-to-lie-down}
\hlchaptermodality{447}

\begin{chapteropening}
\textbf{By the end of this chapter:} \emph{I can complain about a bad night: bedbugs, being bitten, scratching, a sleepless night, and lying down.}

Complete the last lesson of chapter 447: I can complain about a bad night: bedbugs, being bitten, scratching, a sleepless night, and lying down.
\end{chapteropening}

\section[cīmex, cīmicis]{cīmex, cīmicis — a bedbug}

\label{lesson:LA-C447-cimex}



\begin{quote}
Before the new one: say the Latin for a pillow, then the Latin for a blanket.
\end{quote}

\subsection*{You'll want to know: cīmex, cīmicis}
\textbf{cīmex, cīmicis} — \textquotedblleft{}a bedbug\textquotedblright{}.

\begin{quote}
In lectō caupōnae multī cīmicēs sunt.
\end{quote}

\begin{quote}
\emph{In the inn's bed there are lots of bedbugs.}
\end{quote}

Third declension, masculine: \emph{cīmex, cīmicis m.} An old enemy of travellers; also a rude word for a small, annoying man.

\begin{cousinweb}
The origin of \emph{cīmex} is not settled: Lewis \& Short give the sense, a bug, but no etymology, and EtymDB lists it without a parent.
\end{cousinweb}

\begin{grammarlens}[title={What you've built}]
One more word for telling, explaining and describing at B1.
\end{grammarlens}

\subsection*{Your turn}
\begin{itemize}
\raggedright
  \item \emph{Say it:} \emph{cīmex, cīmicis}
  \item \emph{Say it:} \emph{cīmex, cīmicis}, once more
  \item \emph{Say it:} a sentence of your own with \emph{cīmex, cīmicis}
  \item \emph{Recall:} say the Latin for a pillow, then the Latin for a blanket, then say \emph{cīmex, cīmicis} again
\end{itemize}

\begin{tcolorbox}[breakable,colback=teal!4,colframe=teal!35!black,title={Before you move on}]
What does cīmex, cīmicis mean? (\textquotedblleft{}A bedbug\textquotedblright{}.) Say it once more.
\end{tcolorbox}
\section[mordeō, mordēre]{mordeō, mordēre — to bite}

\label{lesson:LA-C447-mordeo}



\begin{quote}
Before the new one: say the Latin for a blanket, then the Latin for a bedbug.
\end{quote}

\subsection*{You'll want to know: mordeō, mordēre}
\textbf{mordeō, mordēre} — \textquotedblleft{}to bite\textquotedblright{}.

\begin{quote}
Tōtam noctem cīmicēs hospitem mordent.
\end{quote}

\begin{quote}
\emph{All night the bedbugs bite the guest.}
\end{quote}

\emph{mordeō, mordēre, momordī, morsum}, with a reduplicated perfect. English \emph{remorse} (a biting back) and \emph{mordant} come from it.

\begin{cousinweb}
\emph{mordeō} is inherited from Proto-Italic \emph{mordeō}, from PIE \emph{(s)merd-} 'to bite, to sting' (EtymDB). English \emph{smart} 'to sting' is a cousin, and \emph{remorse} and \emph{mordant} come from the Latin verb.
\end{cousinweb}

\begin{grammarlens}[title={What you've built}]
One more word for telling, explaining and describing at B1.
\end{grammarlens}

\subsection*{Your turn}
\begin{itemize}
\raggedright
  \item \emph{Say it:} \emph{mordeō, mordēre}
  \item \emph{Say it:} \emph{mordeō, mordēre}, once more
  \item \emph{Say it:} a sentence of your own with \emph{mordeō, mordēre}
  \item \emph{Recall:} say the Latin for a blanket, then the Latin for a bedbug, then say \emph{mordeō, mordēre} again
\end{itemize}

\begin{tcolorbox}[breakable,colback=teal!4,colframe=teal!35!black,title={Before you move on}]
What does mordeō, mordēre mean? (\textquotedblleft{}To bite\textquotedblright{}.) Say it once more.
\end{tcolorbox}
\section[scalpō, scalpere]{scalpō, scalpere — to scratch; to carve}

\label{lesson:LA-C447-scalpo}



\begin{quote}
Before the new one: say the Latin for a bedbug, then the Latin for to bite.
\end{quote}

\subsection*{You'll want to know: scalpō, scalpere}
\textbf{scalpō, scalpere} — \textquotedblleft{}to scratch; to carve\textquotedblright{}.

\begin{quote}
Hospes caput et bracchia scalpit.
\end{quote}

\begin{quote}
\emph{The guest scratches his head and arms.}
\end{quote}

\emph{scalpō, scalpere, scalpsī, scalptum}. Also to carve or engrave. English \emph{scalpel} comes from \emph{scalpellum}, a little knife, from the same root.

\begin{cousinweb}
\emph{scalpō} is inherited from PIE \emph{(s)kel(H)p-}, from a root meaning 'to split, to cut' (EtymDB). English \emph{scalpel} comes from \emph{scalpellum}, a diminutive of \emph{scalprum} 'knife', built on it (Lewis \& Short).
\end{cousinweb}

\begin{grammarlens}[title={What you've built}]
One more word for telling, explaining and describing at B1.
\end{grammarlens}

\subsection*{Your turn}
\begin{itemize}
\raggedright
  \item \emph{Say it:} \emph{scalpō, scalpere}
  \item \emph{Say it:} \emph{scalpō, scalpere}, once more
  \item \emph{Say it:} a sentence of your own with \emph{scalpō, scalpere}
  \item \emph{Recall:} say the Latin for a bedbug, then the Latin for to bite, then say \emph{scalpō, scalpere} again
\end{itemize}

\begin{tcolorbox}[breakable,colback=teal!4,colframe=teal!35!black,title={Before you move on}]
What does scalpō, scalpere mean? (\textquotedblleft{}To scratch; to carve\textquotedblright{}.) Say it once more.
\end{tcolorbox}
\section[īnsomnis, īnsomne]{īnsomnis, īnsomne — sleepless}

\label{lesson:LA-C447-insomnis}



\begin{quote}
Before the new one: say the Latin for to bite, then the Latin for to scratch.
\end{quote}

\subsection*{You'll want to know: īnsomnis, īnsomne}
\textbf{īnsomnis, īnsomne} — \textquotedblleft{}sleepless\textquotedblright{}.

\begin{quote}
Propter cīmicēs viātor īnsomnis in lectō iacet.
\end{quote}

\begin{quote}
\emph{Because of the bedbugs the traveller lies sleepless in bed.}
\end{quote}

An adjective of two endings, from \emph{in-} 'not' and \emph{somnus} 'sleep'. English \emph{insomnia} comes from a noun built on it.

\begin{cousinweb}
\emph{īnsomnis} is built from \emph{in-} 'not' and \emph{somnus} 'sleep' (Lewis \& Short). English \emph{insomnia} comes from the related noun \emph{īnsomnia}.
\end{cousinweb}

\begin{grammarlens}[title={What you've built}]
One more word for telling, explaining and describing at B1.
\end{grammarlens}

\subsection*{Your turn}
\begin{itemize}
\raggedright
  \item \emph{Say it:} \emph{īnsomnis, īnsomne}
  \item \emph{Say it:} \emph{īnsomnis, īnsomne}, once more
  \item \emph{Say it:} a sentence of your own with \emph{īnsomnis, īnsomne}
  \item \emph{Recall:} say the Latin for to bite, then the Latin for to scratch, then say \emph{īnsomnis, īnsomne} again
\end{itemize}

\begin{tcolorbox}[breakable,colback=teal!4,colframe=teal!35!black,title={Before you move on}]
What does īnsomnis, īnsomne mean? (\textquotedblleft{}Sleepless\textquotedblright{}.) Say it once more.
\end{tcolorbox}
\section[cubō, cubāre]{cubō, cubāre — to lie down, to lie (in bed); to sleep}

\label{lesson:LA-C447-cubo}



\begin{quote}
Before the new one: say the Latin for to scratch, then the Latin for sleepless.
\end{quote}

\subsection*{You'll want to know: cubō, cubāre}
\textbf{cubō, cubāre} — \textquotedblleft{}to lie down, to lie (in bed); to sleep\textquotedblright{}.

\begin{quote}
Fessus viātor in lectulō cubat, sed nōn dormit.
\end{quote}

\begin{quote}
\emph{The tired traveller lies on the little bed, but he does not sleep.}
\end{quote}

\emph{cubō, cubāre, cubuī, cubitum}. English \emph{incubate} and \emph{cubicle} come from the same root.

\begin{cousinweb}
\emph{cubō} is inherited from Proto-Italic \emph{kubāō}, from PIE \emph{\'{k}ewb-} 'to lie down' (EtymDB). English \emph{incubate} comes from \emph{incubāre}, built on it (Lewis \& Short).
\end{cousinweb}

\begin{grammarlens}[title={What you've built}]
One more word for telling, explaining and describing at B1.
\end{grammarlens}

\subsection*{Your turn}
\begin{itemize}
\raggedright
  \item \emph{Say it:} \emph{cubō, cubāre}
  \item \emph{Say it:} \emph{cubō, cubāre}, once more
  \item \emph{Say it:} a sentence of your own with \emph{cubō, cubāre}
  \item \emph{Recall:} say the Latin for to scratch, then the Latin for sleepless, then say \emph{cubō, cubāre} again
\end{itemize}

\begin{tcolorbox}[breakable,colback=teal!4,colframe=teal!35!black,title={Before you move on}]
What does cubō, cubāre mean? (\textquotedblleft{}To lie down, to lie (in bed); to sleep\textquotedblright{}.) Say it once more.
\end{tcolorbox}
